\section*{Capítulo \ref{ch:g3:stones-pebbles-sand} --- Piedras, cantos rodados y arena}

\begin{solution}{exo:g3:stones-pebbles-sand:1}
El \omterm{def:g3:stones-pebbles-sand:rock}{lecho rocoso} es la \omterm{def:g3:stones-pebbles-sand:rock}{roca} sólida que hay bajo el suelo en todas partes. Lo
vemos donde no hay suelo que lo esconda: en las laderas empinadas de las
\omterm{def:g1:landscapes-around-us:landform}{montañas}, en los acantilados del \omterm{def:g2:water-on-earth:ocean}{mar}, a lo largo de las carreteras abiertas
a través de una \omterm{def:g1:landscapes-around-us:landform}{colina}.
\end{solution}

\begin{solution}{exo:g3:stones-pebbles-sand:2}
\omterm{def:g3:stones-pebbles-sand:sand}{Barro}, \omterm{def:g3:stones-pebbles-sand:sand}{arena}, \omterm{def:g3:stones-pebbles-sand:pebble}{canto rodado}, bloque.
\end{solution}

\begin{solution}{exo:g3:stones-pebbles-sand:3}
1 milímetro es menos que 2 milímetros: \omterm{def:g3:stones-pebbles-sand:sand}{arena}. 3 milímetros es más que
2 milímetros: \omterm{def:g3:stones-pebbles-sand:pebble}{grava}.
\end{solution}

\begin{solution}{exo:g3:stones-pebbles-sand:4}
Naturales: el \omterm{def:g3:stones-pebbles-sand:pebble}{canto rodado}, el bloque de granito (cortado por la gente, pero
la \omterm{def:g3:stones-pebbles-sand:rock}{roca} es natural), la \omterm{def:g3:stones-pebbles-sand:sand}{arena} de la playa. Hechos por la gente: el ladrillo
y el vidrio.
\end{solution}

\begin{solution}{exo:g3:stones-pebbles-sand:5}
El agua del \omterm{def:g1:landscapes-around-us:water}{río} lo ha hecho rodar por el fondo y lo ha golpeado contra
otras piedras durante mucho tiempo: sus esquinas se han gastado.
\end{solution}

\begin{solution}{exo:g3:stones-pebbles-sand:6}
El granito, para un muro; las losas de piedra finas y planas (una \omterm{def:g3:stones-pebbles-sand:rock}{roca} como
la pizarra), para los tejados.
\end{solution}

\begin{solution}{exo:g3:stones-pebbles-sand:7}
La redonda: cuanto más tiempo hace rodar el agua una piedra, más se gastan
sus esquinas. La afilada no ha viajado lejos de la \omterm{def:g3:stones-pebbles-sand:rock}{roca} de la que se rompió.
\end{solution}

\begin{solution}{exo:g3:stones-pebbles-sand:8}
El \omterm{def:g3:stones-pebbles-sand:pebble}{canto} más grande mide 3 centímetros. Un centímetro son 10 milímetros, y
$10 \div 2 = 5$: caben cinco granos de 2 milímetros a lo largo de
1 centímetro.
\end{solution}

\begin{solution}{exo:g3:stones-pebbles-sand:9}
En la charca tranquila. El agua rápida mantiene en movimiento los granos
diminutos de \omterm{def:g3:stones-pebbles-sand:sand}{barro}; solo donde el agua está casi quieta pueden depositarse
en el fondo.
\end{solution}

\begin{solution}{exo:g3:stones-pebbles-sand:10}
Porque están hechos con materiales sacados del suelo: los ladrillos, de
arcilla; el hormigón, de \omterm{def:g3:stones-pebbles-sand:sand}{arena}, \omterm{def:g3:stones-pebbles-sand:pebble}{grava} y cemento; el vidrio, de \omterm{def:g3:stones-pebbles-sand:sand}{arena}. La
arcilla, la \omterm{def:g3:stones-pebbles-sand:sand}{arena} y la \omterm{def:g3:stones-pebbles-sand:pebble}{grava} vienen de \omterm{def:g3:stones-pebbles-sand:rock}{roca} rota.
\end{solution}
